% -----------------------------------------------------------------------------
% Common mathematical tools for generic QFT, proved before first use
% -----------------------------------------------------------------------------

散射积分反复遇到三个具体困难：能量守恒的 $\delta$ 函数需要在变量变换后给出正确 Jacobian；多个传播子分母需要合并成一个可平移动量的分母；圈动量的张量分子需要利用对称性化成标量积分。这三类工具共同固定散射积分中的 Jacobian、分母组合以及张量约化，并同时给出收敛域与解析延拓条件。

若光滑函数 $f(x)$ 的零点 $x_i$ 都是单根，即 $f(x_i)=0$ 且 $f'(x_i)\neq0$，则复合 $\delta$ 分布满足
\begin{equation}
{
\delta[f(x)]
=\sum_i\frac{\delta(x-x_i)}{|f'(x_i)|}.}
\label{eq:delta-composition-r68}
\end{equation}
这个 Jacobian 因子由变量变换唯一确定，并直接进入质量壳积分与末态相空间测度。

圈积分最常见的问题是多个传播子分母相乘。母公式从每个分母的 Schwinger 参数表示开始。若 $\operatorname{Re}A_i>0$、$\operatorname{Re}\nu_i>0$，则
\begin{equation}
 {
 \frac1{A_i^{\nu_i}}
 =\frac1{\Gamma(\nu_i)}
 \int_0^\infty\dd\tau_i\,\tau_i^{\nu_i-1}\ee^{-A_i\tau_i}.}
 \label{eq:schwinger-parameter-r9}
\end{equation}
令 $\nu\equiv\sum_{i=1}^N\nu_i$，并作变量分解 $\tau_i=sx_i$、$x_i\ge0$、$\sum_i x_i=1$。对公共尺度 $s$ 的积分给出 $\Gamma(\nu)$，于是任意 $N$ 个分母满足
\begin{equation}
{
 \frac1{\prod_{i=1}^N A_i^{\nu_i}}
 =\frac{\Gamma(\nu)}{\prod_{i=1}^N\Gamma(\nu_i)}
 \int_{x_i\ge0}\!\left[\prod_{i=1}^N\dd x_i\,x_i^{\nu_i-1}\right]
 \delta\!\left(1-\sum_{i=1}^N x_i\right)
 \frac1{\left(\sum_{i=1}^N x_iA_i\right)^{\nu}}.}
\label{eq:feynman-n-general-dp74}
\end{equation}
这就是后续圈图使用的母参数化；参数 $x_i$ 无量纲，Minkowski 传播子的 $\ii0^+$ 处方随线性组合一起保留。对 $N=2$，\eqnrefs{eq:feynman-n-general-dp74} 化为
\begin{proposition}[两个分母的 Feynman 参数公式]
先设分母具有正实部，参数实部保证 Schwinger 积分绝对收敛；其余区域按共同的边界值解析延拓。有下列恒等式。
\begin{equation}
{
\frac1{A^{\nu_A}B^{\nu_B}}
=\frac{\Gamma(\nu_A+\nu_B)}{\Gamma(\nu_A)\Gamma(\nu_B)}
\int_0^1\dd x\,
\frac{x^{\nu_A-1}(1-x)^{\nu_B-1}}
{[xA+(1-x)B]^{\nu_A+\nu_B}}.}
\label{eq:feynman-two-general-r68}
\end{equation}
\end{proposition}
\begin{proof}
两个分母的 Feynman 参数并不是独立记忆公式，它直接来自两个 Schwinger 参数的重新参数化。分别对 $A^{\nu_A}$ 和 $B^{\nu_B}$ 使用\eqnrefs{eq:schwinger-parameter-r9}，得到
\begin{align*}
 \frac1{A^{\nu_A}B^{\nu_B}}
 &=\frac1{\Gamma(\nu_A)\Gamma(\nu_B)}
 \int_0^\infty\dd\tau_1\dd\tau_2\,
 \tau_1^{\nu_A-1}\tau_2^{\nu_B-1}
 \ee^{-A\tau_1-B\tau_2}.
\end{align*}
令
\begin{equation*}
 T=\tau_1+\tau_2,
 \qquad x=\frac{\tau_1}{T},
 \qquad
 \tau_1=xT,
 \quad \tau_2=(1-x)T.
\end{equation*}
Jacobian 为 $T$，于是
\begin{align*}
 \frac1{A^{\nu_A}B^{\nu_B}}
 &=\frac1{\Gamma(\nu_A)\Gamma(\nu_B)}
 \int_0^1\dd x\,x^{\nu_A-1}(1-x)^{\nu_B-1}\\
 &\quad\times\int_0^\infty\dd T\,
 T^{\nu_A+\nu_B-1}
 \ee^{-T[xA+(1-x)B]}.
\end{align*}
最后一个积分再次是 $\Gamma$ 函数，直接给出正文\eqnrefs{eq:feynman-two-general-r68}。取 $\nu_A=\nu_B=1$ 即恢复\eqnrefs{eq:feynman-parameter-proof-r9}。
\end{proof}
再取 $\nu_A=\nu_B=1$ 才得到常用特例
\begin{equation}
 {
 \frac1{AB}
 =\int_0^1\dd x\,
 \frac1{[xA+(1-x)B]^2}.}
 \label{eq:feynman-parameter-proof-r9}
\end{equation}
对三个一次分母同样只是 $N=3$、$\nu_i=1$ 的特例：
\begin{proposition}[三个分母的单纯形参数公式]
先设各分母为正实数，指数为正；共用 $\ii0$ 处方的复分母由解析延拓定义。有下列恒等式。
\begin{equation}
 {
 \frac1{ABC}
 =2\int_0^1\dd x\dd y\dd z\,
 \delta(x+y+z-1)
 \frac1{(xA+yB+zC)^3}.}
 \label{eq:feynman-three-denominator-r14}
\end{equation}
\end{proposition}
\begin{proof}
三个分母只比上一个推导多一个比例变量。对 $A,B,C$ 各使用一次幂的 Schwinger 表示，
\begin{equation*}
 \frac1{ABC}
 =\int_0^\infty\dd\tau_1\dd\tau_2\dd\tau_3\,
 \ee^{-A\tau_1-B\tau_2-C\tau_3}.
\end{equation*}
定义
\begin{equation*}
 T=\tau_1+\tau_2+\tau_3,
 \qquad x=\frac{\tau_1}{T},
 \quad y=\frac{\tau_2}{T},
 \quad z=1-x-y,
\end{equation*}
则 $x,y,z\ge0$、$x+y+z=1$，而 Jacobian 为 $T^2$。因此
\begin{align*}
 \frac1{ABC}
 &=\int_0^1\dd x\int_0^{1-x}\dd y
 \int_0^\infty\dd T\,T^2
 \ee^{-T[xA+yB+(1-x-y)C]}\\
 &=2\int_0^1\dd x\int_0^{1-x}\dd y\,
 \frac1{[xA+yB+(1-x-y)C]^3},
\end{align*}
其中用了 $\int_0^\infty T^2\ee^{-DT}\dd T=\Gamma(3)/D^3=2/D^3$。把三角区域写成带 $\delta(x+y+z-1)$ 的对称单纯形测度，就得到正文\eqnrefs{eq:feynman-three-denominator-r14}。
\end{proof}
前因子 $2=\Gamma(3)$ 来自母公式中的公共尺度积分，而不是某张三角圈图的组合因子。

完成多个二次分母的合并和平移积分变量后，许多 $d$ 维欧氏动量积分都会归结为高斯矩。对 $a>0$，$d$ 维欧氏向量 $\ell_E$ 满足
\begin{equation}
 {
 \int\dd^d \ell_E\,\ee^{-a\ell_E^2}
 =\left(\frac\pi a\right)^{d/2}.}
 \label{eq:d-dimensional-gaussian-r9}
\end{equation}
这里的 $d$ 暂时不要求等于物理时空的整数维数。某些量子场论动量积分在 $d=4$ 时发散，却在一段 $\operatorname{Re}d$ 范围内收敛。\emph{维数正规化}把收敛维数范围内得到的解析表达式作为关于 $d$ 的函数，并将它解析延拓到复维数；物理四维由
\begin{equation*}
 d=4-2\epsilon,
 \qquad \epsilon\to0.
\end{equation*}
这样，当积分变量的模 $|\ell|\to\infty$ 时产生的高动量发散会表现为 $1/\epsilon$ 等极点；在量子场论中，这类来自任意大内部动量的发散称为\emph{紫外发散}。有限部分仍由同一个解析函数给出。这里的“非整数维”是给发散积分定义解析延拓的调节参数，并不是说物理时空真的具有分数维数。维数正规化之所以特别适合相对论场论，是因为它可以在计算过程中保持 Lorentz 协变的张量结构；真正的四维物理量要在完成参数重定义和取 $\epsilon\to0$ 极限后得到。

若调节方案保持 $d$ 维 Lorentz 对称性，对只依赖 $\ell^2$ 的标量函数 $f$，二阶张量积分只能与度规成正比：
\begin{equation}
\int\dd^d \ell\,\ell^\mu\ell^\nu f(\ell^2)
=C_f\eta^{\mu\nu}.
\label{eq:rank2-isotropic-ansatz-dp68}
\end{equation}
与 $\eta_{\mu\nu}$ 收缩固定 $C_f$ 后得到
\begin{equation}
 {
 \int\dd^d \ell\,\ell^\mu\ell^\nu f(\ell^2)
 =\frac{\eta^{\mu\nu}}d
 \int\dd^d \ell\,\ell^2 f(\ell^2).}
 \label{eq:tensor-reduction-proof-r9}
\end{equation}
这一步依赖调节器保持平移和 Lorentz 对称性；对尚未调节的发散积分不能先做形式上的“对称平均”。
